% Experiments Section
% For IEEE QCE Submission

\section{Experiments}
\label{sec:experiments}

We evaluate Quantum Mamba and Quantum Hydra on two complementary domains: biomedical time-series classification using two EEG datasets, and natural language understanding using two tasks from the GLUE benchmark.
This pairing tests the same architecture on data with very different statistical properties (continuous multi-channel physiological signals vs.\ discrete tokenized text) and time scales, providing a broad assessment of the proposed models.

% ============================================================
\subsection{EEG datasets}
\label{sec:eeg_datasets}

\subsubsection{PhysioNet Motor Imagery}
\label{sec:physionet}

The PhysioNet EEG Motor Movement/Imagery Database~\cite{schalk2004bci2000,goldberger2000physiobank} is a widely used benchmark for brain--computer interface research. It contains EEG recordings from 109 healthy subjects performing motor imagery tasks, acquired with a 64-channel BCI2000 system at 160\,Hz.
Following standard practice, we use the three runs (runs 4, 8, and 12) in which subjects imagine opening and closing either the \emph{left} or \emph{right} fist, yielding a binary classification task over 64 EEG channels.

\paragraph{Preprocessing.}
Raw EDF recordings are loaded via MNE-Python~\cite{gramfort2013meg}, resampled to a target frequency $f_s \in \{40, 80, 160\}$\,Hz, and epoched to the interval $[1.0, 4.1]$\,s relative to each cue, producing trials of length $T = \lceil 3.1 f_s \rceil$ samples (e.g., $T = 248$ at 80\,Hz).
Signal amplitudes are scaled to millivolts for numerical stability.
Each trial is labeled as left-hand (class 0) or right-hand (class 1) motor imagery.

\paragraph{Splits.}
To assess cross-subject generalization, we perform a \emph{subject-level} stratified split: 70\% of subjects are assigned to the training set, 15\% to validation, and 15\% to test, ensuring no subject appears in more than one split.
We report the mean and standard deviation of test accuracy across three independent splits using seeds $\{2024, 2025, 2026\}$.

\paragraph{Sampling-frequency ablation.}
We train and evaluate each model at $f_s = 40$, $80$, and $160$\,Hz in order to study how quantum advantage scales with temporal resolution and sequence length.
This yields three sequence lengths ($T = 124, 248, 496$) with identical content but different quantum circuit context windows, allowing us to isolate the effect of sequence length from confounding factors.

\subsubsection{ADFTD Dementia EEG}
\label{sec:adftd}

The ADFTD dataset~\cite{miltiadous2023dataset}\footnote{OpenNeuro accession ds004504.} is a publicly available resting-state, eyes-closed EEG corpus collected at the 2nd Department of Neurology of AHEPA General Hospital (Thessaloniki, Greece) for the study of dementia.
It contains recordings from 88 subjects: 36 with Alzheimer's disease (AD), 23 with frontotemporal dementia (FTD), and 29 healthy controls (CN).
Data were acquired with a Nihon Kohden EEG 2100 clinical device using 19 scalp electrodes placed according to the international 10--20 system (Fp1, Fp2, F3, F4, C3, C4, P3, P4, O1, O2, F7, F8, T3, T4, T5, T6, Fz, Cz, Pz) at a sampling rate of 500\,Hz.
Recording duration varied from 5 to 21 minutes per subject.
We use the pre-processed \texttt{derivatives/} version released with the dataset, which has already been (i)~band-pass filtered (0.5--45\,Hz, Butterworth), (ii)~re-referenced to A1--A2, (iii)~artifact-corrected using the Artifact Subspace Reconstruction (ASR) algorithm, and (iv)~cleaned with ICA-based removal of eye and jaw components via ICLabel.

\paragraph{Preprocessing.}
For each subject, we load the .set file using MNE-Python and convert the signal to microvolts.
The continuous recording is then partitioned into non-overlapping fixed-length segments of 2{,}500 samples (5\,seconds at 500\,Hz).
This segmentation yields a consistent trial length across subjects and provides enough per-subject data for batch-level training.
We apply per-channel z-score normalization using statistics computed only from the training set, and clip values to $\pm 10\sigma$ to suppress residual outliers.
Each segment inherits the diagnostic label of its source subject.

\paragraph{Task setup and splits.}
We consider the three-class classification task (AD vs.\ FTD vs.\ CN) over the 19-channel $\times$ 2{,}500-sample trials.
As with PhysioNet, we use a \emph{subject-level} stratified split (70\% train / 15\% val / 15\% test) to prevent leakage of subject-specific features between splits, and we stratify by diagnostic group so that the class distribution is preserved across splits.
Results are averaged over three seeds.

% ============================================================
\subsection{GLUE benchmark}
\label{sec:glue}

The General Language Understanding Evaluation (GLUE) benchmark~\cite{wang2018glue} is a standard suite of nine English-language natural language understanding tasks used to evaluate the transferability and robustness of sentence-level models.
From this suite, we select two contrasting tasks---SST-2 and RTE---that together stress both abundant-data single-sentence modeling and data-scarce sentence-pair reasoning.

\subsubsection{SST-2: Stanford Sentiment Treebank (binary)}
\label{sec:sst2}

SST-2~\cite{socher2013sst} is a binary sentiment classification task built from movie reviews from Rotten Tomatoes.
Each example is a single sentence labeled as either \emph{positive} or \emph{negative}.
The GLUE release provides approximately 67{,}000 training examples and 872 validation examples, making it by far the largest of the two GLUE tasks we consider and a natural testbed for the models' raw classification capacity.

\subsubsection{RTE: Recognizing Textual Entailment (binary)}
\label{sec:rte}

RTE~\cite{dagan2005pascal,bentivogli2009fifth} is a sentence-pair entailment task aggregated from a series of textual entailment challenges.
Each example consists of a \emph{premise} and a \emph{hypothesis}, and the model must predict whether the premise \emph{entails} the hypothesis (label 0) or not (label 1).
The GLUE release provides only 2{,}490 training examples and 277 validation examples, making RTE a data-scarce, low-resource regime that probes generalization from limited supervision---a setting in which inductive biases of the model architecture become especially important.

\subsubsection{Preprocessing and model interface}
\label{sec:glue_preproc}

Both tasks are loaded directly from the HuggingFace \texttt{datasets} library.
Inputs are tokenized with the pre-trained \texttt{bert-base-uncased} WordPiece tokenizer~\cite{devlin2019bert}: SST-2 sentences are passed as single-segment inputs, while RTE premise--hypothesis pairs are joined using the standard \texttt{[CLS]\,premise\,[SEP]\,hypothesis\,[SEP]} format.
All sequences are padded or truncated to a fixed length of $T = 128$ tokens, with an attention mask tracking valid positions.
Token IDs are then passed through a trainable embedding layer of dimension $d$ to produce a continuous sequence $\mathbf{X} \in \mathbb{R}^{B \times 128 \times d}$, which serves as the input to Quantum Mamba or Quantum Hydra in place of the raw feature matrix used for EEG (Eq.~\eqref{eq:input_proj}).
In this setting, the ``channel'' dimension $F$ of our models corresponds to the token embedding dimension, and the chunked temporal processing of Section~\ref{sec:qm_core} operates over token positions instead of time steps.

\paragraph{Labels and metrics.}
Both SST-2 and RTE are binary classification tasks with a single \texttt{label} field; we train with standard cross-entropy loss and report accuracy on the GLUE development set, following the common GLUE evaluation protocol since test-set labels are hidden behind an online submission server.
As with the EEG experiments, results are averaged over three random seeds.
